\documentclass[../report.tex]{subfiles}
\begin{document}

    \begin{frame}
        \frametitle{3a: Haar Features}
        \begin{figure}[!htb]
            \centering
            \frame{\includegraphics[keepaspectratio,height=0.65\textheight,width=0.45\textwidth]{ps6-3-a-1}}
            \caption{ps6-3-a-1}
        \end{figure}
    \end{frame}

    \begin{frame}
        \frametitle{3a: Haar Features}
        \begin{figure}[!htb]
            \centering
            \frame{\includegraphics[keepaspectratio,height=0.65\textheight,width=0.45\textwidth]{ps6-3-a-2}}
            \caption{ps6-3-a-2}
        \end{figure}
    \end{frame}

    \begin{frame}
        \frametitle{3a: Haar Features}
        \begin{figure}[!htb]
            \centering
            \frame{\includegraphics[keepaspectratio,height=0.65\textheight,width=0.45\textwidth]{ps6-3-a-3}}
            \caption{ps6-3-a-3}
        \end{figure}
    \end{frame}

    \begin{frame}
        \frametitle{3a: Haar Features}
        \begin{figure}[!htb]
            \centering
            \frame{\includegraphics[keepaspectratio,height=0.65\textheight,width=0.45\textwidth]{ps6-3-a-4}}
            \caption{ps6-3-a-4}
        \end{figure}
    \end{frame}

    \begin{frame}
        \frametitle{3a: Haar Features}
        \begin{figure}[!htb]
            \centering
            \frame{\includegraphics[keepaspectratio,height=0.65\textheight,width=0.45\textwidth]{ps6-3-a-5}}
            \caption{ps6-3-a-5}
        \end{figure}
    \end{frame}

    \begin{frame}[t]
        \frametitle{3c: Analysis}
        \begin{normalsize}
            \begin{itemize}
                \setlength\itemsep{1em}\fontsize{6pt}{6pt}

                \item[]{\textbf{\selectfont\textcolor{blue}{ How does working with integral images help with computation time? Give some examples comparing this method and np.sum. }}}
                
                \item[]\textbf{There are many advantages in using integral images.\\
Once the integral image is computed from an image.\\
Obtaining a rectangular summation in any section of  image is in\\
near constant time as there is only one addition and\\
two subtraction operators in A - B - C + D\\
np.sum() is computed in linear time, and computing the sum of an image\\
at a particular region is O(N^2) time.\\
Using np.sum would not be a scalable solution, especially in adaboost,\\
an algorithm where  multiple sections of the image have to be summed.\\
}
            \end{itemize}
        \end{normalsize}
    \end{frame}
    
\end{document}